% !TeX root = ../complete.tex
\section{مصرف انرژی در سطح معماری}
\label{sec:power-architecture}

\subsection{محل انجام کار، توازی و \lr{Pipeline}}
\label{sec:power-parallelism-pipeline}

تا اینجا انتخاب‌هایی مانند تغییر اندازه گیت، فعالیت و ولتاژ را
برای یک مدار بررسی کردیم. در سطح \lr{Microarchitecture}، می‌توان
تعیین کرد کار بین چه واحدهایی تقسیم شود و داده چه مسیری طی کند.
قرار دادن واحدهای پرارتباط نزدیک یکدیگر می‌تواند ظرفیت سیم‌ها را
کم کند. یک واحد ویژه برای کاری پرتکرار ممکن است با انتقال داده
و فعالیت منطقیِ کم‌تر همان نتیجه را تولید کند. مقایسه‌ی معماری‌ها
به نوع کار، مسیر حرکت داده و انرژی اجزای کنترلی وابسته است
\cite[بخش \secnum{۵.۵.۱}، صفحه ۲۰۴]{weste}.

\lr{Latency} زمان شروع یک کار تا آماده شدن نتیجه‌ی آن است.
\lr{Throughput} تعداد نتیجه‌های کامل‌شده در واحد زمان است. فرض
کنید یک واحد در هر دوره ساعت یک نتیجه آماده می‌کند. اگر فرکانس آن
\(f\) باشد، در حالت پایدار نرخ خروجی \(f\) نتیجه بر ثانیه است.
در معماری موازی، \(N\) نسخه‌ی همان واحد روی کارهای مستقل کار
می‌کنند. اگر هر نسخه با فرکانس \(f/N\) یک نتیجه در هر دوره
بدهد، نرخ کل برابر \(N(f/N)=f\) است. در این مدل، تأخیر یک کار
با دوره‌ی طولانی‌تر هر واحد افزایش می‌یابد. فرکانس کم‌تر می‌تواند
اجازه‌ی کاهش ولتاژ را بدهد، به شرط آنکه واحد در دوره‌ی جدید پاسخ
خود را کامل کند \cite[بخش \secnum{۵.۵.۲} و شکل
\secnum{۵.۲۹}، صفحات ۲۰۴ و ۲۰۵]{weste}.

اگر یک نتیجه به همان تعداد شارژ در یکی از واحدها نیاز داشته باشد
و ولتاژ کار به ۸۰ درصد مقدار اولیه برسد، انرژی سوئیچینگِ محاسبه هر
نتیجه در آن واحد به نسبت \(0.8^2=0.64\) کاهش می‌یابد. برای انرژی
کل باید مصرف واحدهای اضافی، نشتی آن‌ها، شبکه ساعت و مدار انتخاب
خروجی را هم حساب کرد. این نکته به‌ویژه هنگامی مهم است که نرخ ورود
کار برای استفاده از همه واحدها کافی نباشد.

در \lr{Pipeline}، کار به چند مرحله تقسیم می‌شود و میان مراحل
\lr{Register} قرار می‌گیرد. در هر چرخه مرحله اول می‌تواند کار
جدیدی را آغاز کند، درحالی‌که مراحل بعد روی ورودی‌های قبلی کار
می‌کنند. پس از پر شدن مراحل، اگر هر مرحله در یک دوره پاسخ دهد
و توقفی رخ ندهد، در هر دوره یک نتیجه آماده می‌شود. برای مثال،
دو بخش با تأخیرهای \qty{4}{\nano\second} و
\qty{4}{\nano\second} را در نظر بگیرید. در مدل فرضی، مجموع
هزینه زمانی مربوط به \lr{Register} در هر مرحله
\qty{1}{\nano\second} است. با یک مرحله، دوره ساعت دست‌کم
\qty{9}{\nano\second} است. با دو مرحله، دوره می‌تواند به
\qty{5}{\nano\second} برسد و یک کار پس از دو دوره، یعنی
\qty{10}{\nano\second} کامل شود. در نتیجه، پس از پر شدن
\lr{Pipeline} فاصله خروجی‌ها کم‌تر شده و زمان رسیدن هر کار به
خروجی کمی افزایش یافته است. اگر هدف حفظ نرخ خروجی اولیه باشد،
از این فرصت زمانی می‌توان برای کاهش ولتاژ استفاده کرد. سربار
ظرفیت و ساعت \lr{Register}ها و توازن تأخیر مراحل باید در برآورد
انرژی وارد شوند \cite[بخش \secnum{۵.۵.۲}، صفحه ۲۰۵]{weste}.

\subsection{حالت‌های کار و کنترل برآوردها}
\label{sec:power-management-modes}

در حالت فعال، بلوک کار انجام می‌دهد. در حالت بی‌کار، می‌توان با
\lr{Clock Gating} گذارها را کم کرد یا بسته به بار کاری، ولتاژ و
فرکانس را تغییر داد. برای خواب عمیق‌تر، \lr{Power Gating} نشتی
بخشی از مدار را کاهش می‌دهد و هزینه بیداری دارد. انتخاب حالت به
زمان پیش‌بینی‌شده بی‌کاری، زمان سربه‌سر انرژی در رابطه
\eqref{eq:power-gating-break-even} و مهلت پاسخ بعدی وابسته است.
زمان بی‌کاری کوتاه ممکن است فقط فرصت توقف ساعت را بدهد. در بازه
بلندتر، صرفه‌جویی نشتی می‌تواند انرژی تغییر حالت را جبران کند
\cite[بخش \secnum{۵.۵.۳}، صفحات ۲۰۵ و ۲۰۶]{weste}.

برای نمونه، تغییر هدف از کمینه کردن تأخیر به انجام کار در یک
مهلت مشخص، گاهی اجازه می‌دهد گیت‌ها کوچک‌تر شوند. در مقابل، اگر
در مدل \lr{Logical Effort} بار هر مرحله هم با گیت راه‌اندازش
به یک نسبت رشد کند، نسبت بار به ظرفیت ورودی ثابت می‌ماند. در این
شرایط، رشد همگانی اندازه می‌تواند ظرفیت قابل شارژ را بیشتر کند،
درحالی‌که مدل برای آن مراحل بهبود تأخیر پیش‌بینی نمی‌کند. وضعیت
بار بیرونی و ظرفیت سیم هنگام استفاده از این استدلال باید مشخص باشد
\cite[بخش \secnum{۵.۶}، صفحه ۲۰۶]{weste}.

در گزارش نتیجه، توان را همراه با بار کاری، ولتاژ و نرخ کار بیان
کنید. برای مقایسه انرژی به‌ازای کار، از نرخ تکمیل کار در بخش
\ref{sec:power-energy-per-operation} استفاده کنید.
\lr{Power--Delay Product} یا \lr{PDP}، در یک عملیات حاصل‌ضرب توان
متوسط در مدت همان عملیات است و واحد ژول دارد. \lr{EDP} حاصل‌ضرب
انرژی عملیات در تأخیر است و واحد ژول‌ثانیه دارد. همچنین سهم
نشتی را هم در کار فعال و هم در زمان انتظار در نظر بگیرید، چون
طول هر دو بازه به نحوه کار سامانه وابسته است
\cite[بخش \secnum{۵.۶}، صفحه ۲۰۶]{weste}.

دید تاریخی فصل \lr{Power} نیز همین تغییر موازنه را دنبال می‌کند.
در منطق قدیمی \lr{nMOS} با بار همیشه روشن، بعضی خروجی‌های ثابت
مسیر ماندگار جریان داشتند. فراگیر شدن \lr{CMOS} مکمل این سهم را
کاهش داد. در فناوری‌های بعدی، کاهش آستانه و نازک شدن دی‌الکتریک
گیت باعث شد کنترل نشتی، در کنار کاهش ولتاژ و فعالیت، دوباره اهمیت
پیدا کند. این روند توضیح می‌دهد چرا برای برآورد توان باید مدل
افزاره و معماریِ همان طراحی را با هم دید
\cite[بخش \secnum{۵.۷}، صفحات ۲۰۷ و ۲۰۸]{weste}.
